% !TEX root = ../main.tex

\chapter{Szenarien und Use Cases}

Um die funktionalen Anforderungen an das Assistenzsystem \glqq Ascenta\grqq\ zu definieren und die spätere technische Umsetzung zu rechtfertigen, werden im Folgenden konkrete Anwendungsszenarien (Use Cases) beschrieben. Diese Szenarien repräsentieren alltägliche Herausforderungen von stark sehbeeinträchtigten oder blinden Personen, die durch das System gelöst werden. Gemäß dem aktuellen Entwicklungsstand fokussieren sich die Anwendungsfälle auf die Erkennung von Personen im direkten Umfeld, die Distanzabschätzung sowie die auditive Interaktion.

\section{Use Case A: Kollisionsvermeidung mit Personen}

\subsection{Ausgangssituation}
Ein Nutzer bewegt sich mit einem regulären Langstock durch einen öffentlichen Raum, beispielsweise ein Foyer oder einen Gehweg. Während der Langstock Bodenhindernisse (wie Treppenstufen oder Bordsteinkanten) zuverlässig ertastet, registriert er keine stillstehenden oder sich nähernden Personen auf Körper- oder Kopfhöhe, bevor es zu einem physischen Kontakt kommt.

\subsection{Systemverhalten}
Das Kamerasystem der Ascenta-Brille erfasst kontinuierlich das Sichtfeld des Nutzers. Sobald eine Person in den Erfassungsbereich tritt, wird diese durch das lokale Machine-Learning-Modell (FOMO) auf dem Mikrocontroller identifiziert. Simultan misst der Time-of-Flight-Sensor (ToF) die exakte Distanz zu diesem Objekt. 

Unterschreitet die Person eine vordefinierte kritische Distanz (beispielsweise einen Meter), fusioniert die Smartphone-Applikation diese Sensordaten und generiert unmittelbar ein akustisches Warnsignal. Der Nutzer wird dadurch in Echtzeit auf das Hindernis aufmerksam gemacht und kann rechtzeitig ausweichen, ohne dass ein physischer Kontakt entsteht.

\section{Use Case B: Kontextbezogene Umgebungsabfrage per Sprachbefehl}

\subsection{Ausgangssituation}
Der Nutzer befindet sich in einer Warteschlange oder an einem Haltepunkt (z. B. einer Bushaltestelle) und möchte wissen, ob sich eine Person direkt vor ihm befindet, um entsprechend aufrücken oder Abstand halten zu können. Ein weites Ausschwenken des Langstocks ist in dieser Situation unerwünscht, da es andere Personen stören oder irritieren könnte.

\subsection{Systemverhalten}
Anstatt die Umgebung physisch abzutasten, nutzt der Anwender die integrierte Sprachsteuerung. Der Nutzer aktiviert das System (beispielsweise durch ein definiertes Interaktionsmuster an der Brille) und stellt verbal eine Frage, wie etwa: \glqq Befindet sich jemand vor mir?\grqq. 

Das Audiosignal wird über das in der Brille verbaute Mikrofon erfasst und via UDP-Protokoll latenzarm an das verbundene Smartphone gestreamt. Dort verarbeitet die lokale Vosk-Engine den Sprachbefehl offline, um den Datenschutz zu gewährleisten. Die Applikation gleicht die Anfrage mit den aktuellen Kameradaten (Personenerkennung) und der Distanzmessung (ToF) ab. Anschließend erhält der Nutzer über die Sprachausgabe (Text-to-Speech) ein diskretes, akustisches Feedback, beispielsweise: \glqq Person in einem Punkt fünf Metern Entfernung erkannt.\grqq\ Dadurch gewinnt der Nutzer ohne stigmatisierende Bewegungen ein präzises Verständnis seiner unmittelbaren sozialen Umgebung.